[# --- Herbruikbare TikZ-grafiekmacro's (Jinja) --- #]
[% macro hbarchart(items, unit) %]
\begin{center}
\begin{tikzpicture}[y=-0.72cm]
[% set vmax = (items|map(attribute='value')|max) or 1 %]
[% for it in items %]
  \node[left,font=\small] at (0,[[ loop.index0 ]]) {[[ it.label | latex_escape ]]};
  \fill[{[[ it.color ]]}] (0.15,[[ loop.index0 ]]-0.26) rectangle ({0.15+[[ it.value ]]/[[ vmax ]]*6},[[ loop.index0 ]]+0.26);
  \node[right,font=\small] at ({0.2+[[ it.value ]]/[[ vmax ]]*6},[[ loop.index0 ]]) {[[ it.value ]][% if unit %] [[ unit ]][% endif %]};
[% endfor %]
\end{tikzpicture}
\end{center}
[% endmacro %]
[% macro stackbar(items) %]
[% set tot = (items|sum(attribute='value')) or 1 %]
[% set ns = namespace(acc=0.0) %]
\resizebox{\linewidth}{!}{%
\begin{tikzpicture}[x=1cm,y=1cm]
[% for it in items %]
[% if it.value > 0 %]
  \fill[{[[ it.color ]]}] ([[ ns.acc/tot*12 ]],0) rectangle ([[ (ns.acc+it.value)/tot*12 ]],0.7);
  \node[white,font=\scriptsize] at ([[ (ns.acc+it.value*0.5)/tot*12 ]],0.35) {[[ it.value ]]};
[% set ns.acc = ns.acc + it.value %]
[% endif %]
[% endfor %]
  \draw (0,0) rectangle (12,0.7);
\end{tikzpicture}}
\par\vspace{2pt}
{\scriptsize [% for it in items %][% if it.value > 0 %]\textcolor{[[ it.color ]]}{\rule{1.4ex}{1.4ex}}~[[ it.label | latex_escape ]] ([[ it.value ]])\quad[% endif %][% endfor %]}
[% endmacro %]
[% macro piechart(slices) %]
\resizebox{\linewidth}{!}{%
\begin{tikzpicture}
[% for s in slices %]
  \fill[{[[ s.color ]]}] (0,0) -- ([[ s.a0 ]]:2) arc ([[ s.a0 ]]:[[ s.a1 ]]:2) -- cycle;
[% endfor %]
  \fill[white] (0,0) circle (1);
[% for s in slices %]
  \draw[white,line width=1.2pt] (0,0) -- ([[ s.a0 ]]:2);
[% endfor %]
\end{tikzpicture}}
\par\vspace{2pt}
{\scriptsize [% for s in slices %]\textcolor{[[ s.color ]]}{\rule{1.4ex}{1.4ex}}~[[ s.label | latex_escape ]] ([[ s.count ]])\quad[% endfor %]}
[% endmacro %]
[% macro histchart(hist, label) %]
\resizebox{\linewidth}{!}{%
\begin{tikzpicture}[x=0.9cm,y=1cm]
[% set hmax = hist.max or 1 %]
[% for b in hist.bins %]
[% if b.count > 0 %]
  \fill[cblue] ([[ b.x + 0.12 ]],0) rectangle ([[ b.x + 0.88 ]],{[[ b.count ]]/[[ hmax ]]*3.0});
  \node[above,font=\tiny] at ([[ b.x + 0.5 ]],{[[ b.count ]]/[[ hmax ]]*3.0}) {[[ b.count ]]};
[% endif %]
  \node[below,font=\tiny] at ([[ b.x + 0.5 ]],0) {[[ b.x ]]};
[% endfor %]
  \draw[gray] (0,0) -- (10,0);
  \node[below=1.2em,font=\small] at (5,0) {[[ label | latex_escape ]]};
\end{tikzpicture}}
[% endmacro %]
[% macro slagantt(rows) %]
\begin{center}
\begin{tikzpicture}[y=-0.52cm,x=1cm]
[% for r in rows %]
  \node[left,font=\scriptsize] at (0,[[ loop.index0 ]]) {[[ r.id ]]};
  \fill[{[[ r.color ]]}] (0.1,[[ loop.index0 ]]-0.2) rectangle ({0.1+ln([[ r.days ]]+1)/ln(366)*10},[[ loop.index0 ]]+0.2);
  \node[right,font=\scriptsize] at ({0.2+ln([[ r.days ]]+1)/ln(366)*10},[[ loop.index0 ]]) {[[ r.sla | latex_escape ]]};
[% endfor %]
\end{tikzpicture}
\end{center}
[% endmacro %]

% --- Colofon ---
\clearpage
[% if taal == 'en' %]
\section*{Colophon}
[% else %]
\section*{Colofon}
[% endif %]
\renewcommand{\arraystretch}{1.35}
\begin{tabular}{@{}p{4.2cm}p{10cm}@{}}
  [% if taal == 'en' %]\textbf{Title}[% else %]\textbf{Titel}[% endif %] & \rapporttitel \\
  [% if subtitel %][% if taal == 'en' %]\textbf{Subtitle}[% else %]\textbf{Subtitel}[% endif %] & \rapportsubtitel \\[% endif %]
  [% if rapport_versie %][% if taal == 'en' %]\textbf{Version}[% else %]\textbf{Versie}[% endif %] & [[ rapport_versie ]] \\[% endif %]
  [% if taal == 'en' %]\textbf{Status}[% else %]\textbf{Status}[% endif %] & [% if rapport_versie_status == 'definitief' or rapport_versie_status == 'final' %][% if taal == 'en' %]Final[% else %]Definitief[% endif %][% elif rapport_versie_status == 'review' %][% if taal == 'en' %]Under review[% else %]Ter review[% endif %][% else %][% if taal == 'en' %]Draft[% else %]Concept[% endif %][% endif %] \\
  [% if taal == 'en' %]\textbf{Date}[% else %]\textbf{Datum}[% endif %] & \rapportdatum \\
  [% if taal == 'en' %]\textbf{Classification}[% else %]\textbf{Classificatie}[% endif %] & \textcolor{\rapporttlpkleur}{\rapportclassificatie} \\
  [% if project %][% if taal == 'en' %]\textbf{Client}[% else %]\textbf{Opdrachtgever}[% endif %] & \rapportproject \\[% endif %]
  [% if taal == 'en' %]\textbf{Author}[% else %]\textbf{Opsteller}[% endif %] & \rapportauteur \\
  [% if taal == 'en' %]\textbf{Assessment type}[% else %]\textbf{Type onderzoek}[% endif %] & [% if testtype == 'redteam' %]Red team assessment[% else %][% if taal == 'en' %]Penetration test[% else %]Penetratietest[% endif %][% endif %][% if testscope %] ([[ testscope | latex_escape ]])[% endif %] \\
  [% if rapport_omgeving %][% if taal == 'en' %]\textbf{Environment}[% else %]\textbf{Omgeving}[% endif %] & [[ rapport_omgeving ]] \\[% endif %]
  [% if test_start or test_end %][% if taal == 'en' %]\textbf{Test period}[% else %]\textbf{Testperiode}[% endif %] & [[ test_start or '?' ]] -- [[ test_end or '?' ]] \\[% endif %]
  [% if rapport_contact %]\textbf{Contact} & [[ rapport_contact ]] \\[% endif %]
\end{tabular}
\renewcommand{\arraystretch}{1.0}

\vspace{1.2cm}
\noindent{\itshape\small
[% if taal == 'en' %]
This document is classified \rapportclassificatie{} and is intended solely for the client named above. Distribution to third parties is only permitted with prior written consent.
[% else %]
Dit document is geclassificeerd als \rapportclassificatie{} en is uitsluitend bedoeld voor de hierboven genoemde opdrachtgever. Verspreiding aan derden is alleen toegestaan na voorafgaande schriftelijke toestemming.
[% endif %]
}

[% if rapport_revisies %]
\vspace{1cm}
[% if taal == 'en' %]
\subsection*{Revision history}
[% else %]
\subsection*{Revisiehistorie}
[% endif %]
\renewcommand{\arraystretch}{1.3}
\begin{tabular}{@{}p{1.8cm}p{2.6cm}p{3.2cm}p{5.6cm}@{}}
  \rowcolor{rijkblauw}[% if taal == 'en' %]\textcolor{white}{\textbf{Version}} & \textcolor{white}{\textbf{Date}} & \textcolor{white}{\textbf{Author}} & \textcolor{white}{\textbf{Change}}[% else %]\textcolor{white}{\textbf{Versie}} & \textcolor{white}{\textbf{Datum}} & \textcolor{white}{\textbf{Auteur}} & \textcolor{white}{\textbf{Wijziging}}[% endif %] \\
  \hline
[% for rv in rapport_revisies %]
  [[ rv.versie | latex_escape ]] & [[ rv.datum | latex_escape ]] & [[ rv.auteur | latex_escape ]] & [[ rv.wijziging | latex_escape ]] \\
[% endfor %]
\end{tabular}
\renewcommand{\arraystretch}{1.0}
[% endif %]

[% if rapport_distributie %]
\vspace{1cm}
[% if taal == 'en' %]
\subsection*{Distribution}
[% else %]
\subsection*{Distributie}
[% endif %]
\renewcommand{\arraystretch}{1.3}
\begin{tabular}{@{}p{3.6cm}p{3.6cm}p{3.6cm}p{2.4cm}@{}}
  \rowcolor{rijkblauw}[% if taal == 'en' %]\textcolor{white}{\textbf{Name}} & \textcolor{white}{\textbf{Organisation}} & \textcolor{white}{\textbf{Role}} & \textcolor{white}{\textbf{Version}}[% else %]\textcolor{white}{\textbf{Naam}} & \textcolor{white}{\textbf{Organisatie}} & \textcolor{white}{\textbf{Rol}} & \textcolor{white}{\textbf{Versie}}[% endif %] \\
  \hline
[% for d in rapport_distributie %]
  [[ d.naam | latex_escape ]] & [[ d.organisatie | latex_escape ]] & [[ d.rol | latex_escape ]] & [[ d.versie | latex_escape ]] \\
[% endfor %]
\end{tabular}
\renewcommand{\arraystretch}{1.0}
[% endif %]
\clearpage

[% if sections.managementsamenvatting and management_samenvatting %]
% --- Managementsamenvatting (vóór de inhoudsopgave; vaste positie, buiten reorder) ---
[% if taal == 'en' %]
\section{Executive Summary}
[% else %]
\section{Managementsamenvatting}
[% endif %]

[[ management_samenvatting ]]

[% if algeheel_beeld %]
[% if taal == 'en' %]\subsection*{Overall assessment}[% else %]\subsection*{Algeheel beeld}[% endif %]
[[ algeheel_beeld ]]
[% endif %]

[% if sections.top_risk_gauge and top_risk_gauge %]
[% if taal == 'en' %]\subsection*{Highest risk}[% else %]\subsection*{Hoogste risico}[% endif %]
\ibgauge{[[ top_risk_gauge['gauge_angle'] ]]}{[[ top_risk_gauge['score'] ]]}{[[ top_risk_gauge['label'] | latex_escape ]]}{[% if taal == 'en' %]Highest CVSS base score[% else %]Hoogste CVSS-basescore[% endif %][% if top_risk_gauge['nr'] %] --- [[ top_risk_gauge['nr'] ]] [[ top_risk_gauge['titel'] | latex_escape ]][% endif %]}
[% endif %]

[[ key_findings_tex ]]

[[ digid_knockout_tex ]]
\clearpage
[% endif %]

% --- Inhoudsopgave ---
\tableofcontents
% Lijst van figuren, alleen als er bewijsbeelden zijn (anders geen lege kop).
[% if findings | selectattr('evidence') | list %]
\newpage
\listoffigures
[% endif %]
\newpage

% --- Technische rapportage (hoofdstuk met de technische items als subsecties) ---
[% if taal == 'en' %]
\section{Technical report}
[% else %]
\section{Technische rapportage}
[% endif %]

% --- Inleiding ---
[% if taal == 'en' %]
\subsection{Introduction}
[% else %]
\subsection{Inleiding}
[% endif %]

[% if rapport_inleiding %]
[[ rapport_inleiding ]]
[% else %]
[% if testtype == 'redteam' %]
[% if taal == 'en' %]
This report contains the findings of the red team assessment performed on \rapportproject.
The objective was to simulate a realistic attack to evaluate the detection and response capabilities
of the organisation, as well as to identify vulnerabilities in people, processes and technology.
[% else %]
Dit rapport bevat de bevindingen van het red team assessment uitgevoerd op \rapportproject.
Het doel was het simuleren van een realistische aanval om de detectie- en responscapaciteiten
van de organisatie te evalueren, evenals het identificeren van kwetsbaarheden in mensen,
processen en technologie.
[% endif %]
[% else %]
[% if taal == 'en' %]
This report contains the findings of the penetration test performed on \rapportproject.
The objective was to identify vulnerabilities that could compromise the confidentiality,
integrity and availability of the systems.
[% else %]
Dit rapport bevat de bevindingen van de penetratietest uitgevoerd op \rapportproject.
Het doel van de test was het identificeren van kwetsbaarheden die de vertrouwelijkheid,
integriteit en beschikbaarheid van de systemen kunnen aantasten.
[% endif %]
[% endif %]
[% endif %]

[% if taal == 'en' %]
The test was performed by \rapportauteur{} on \rapportdatum.
[% else %]
De test is uitgevoerd door \rapportauteur{} op \rapportdatum.
[% endif %]

\subsubsection{Scope}

[# Scope-targets zoals ingesteld in dashboard/report; anders terugval op de
   locaties uit de bevindingen. #]
[% if scope_targets %]
[% if taal == 'en' %]
The following targets are within scope:
[% else %]
De volgende targets vallen binnen de scope:
[% endif %]

{\arrayrulecolor{white}\renewcommand{\arraystretch}{1.25}%
\begin{xltabular}{\textwidth}{|l|l|X|}
  \hline
  \rowcolor{rijkblauw}\textcolor{white}{\textbf{Target}} & \textcolor{white}{\textbf{Type}} & \textcolor{white}{\textbf{[% if taal == 'en' %]Description[% else %]Beschrijving[% endif %]}} \\ \hline
  \endhead
[% for st in scope_targets %]
  [[ (st.target or '') | latex_escape ]] & [[ (st.type or '') | latex_escape ]] & [[ (st.beschrijving or '') | latex_escape ]] \\ \hline
[% endfor %]
\end{xltabular}}
[% else %]
[% if taal == 'en' %]
The following systems were part of the test:
[% else %]
De volgende systemen zijn onderdeel van de test:
[% endif %]
\begin{itemize}
[% for loc in scoped_hosts %]
  \item [[ loc | latex_escape ]]
[% endfor %]
[% if not scoped_hosts %]
[% if taal == 'en' %]
  \item See findings for specific locations
[% else %]
  \item Zie bevindingen voor specifieke locaties
[% endif %]
[% endif %]
\end{itemize}
[% endif %]

[# Kritieke functies / kroonjuwelen (ART), bij de scope. #]
[% if kritieke_functies %]
\subsubsection{[% if taal == 'en' %]Critical functions[% else %]Kritieke functies[% endif %]}
[[ kritieke_functies ]]
[% endif %]

[% if sections.methodologie %]
%SEC:methodologie%
% --- Methodologie ---
[% if taal == 'en' %]
\subsection{Methodology}
[% else %]
\subsection{Methodologie}
[% endif %]

[% if rapport_methodologie %]
[[ rapport_methodologie ]]
[% else %]
[% if testtype == 'redteam' %]
[% if taal == 'en' %]
The red team assessment was performed in accordance with the MITRE ATT\&CK framework and the Cyber Kill Chain.
The test consisted of the following phases:

\begin{enumerate}
  \item \textbf{Reconnaissance} -- OSINT and target identification
  \item \textbf{Weaponisation} -- Prepare tooling and payloads
  \item \textbf{Delivery} -- Obtain initial access
  \item \textbf{Exploitation} -- Exploit vulnerabilities
  \item \textbf{Installation} -- Establish persistence
  \item \textbf{Command \& Control} -- Set up C2 channel
  \item \textbf{Actions on Objectives} -- Achieve goals
\end{enumerate}
[% else %]
Het red team assessment is uitgevoerd conform het MITRE ATT\&CK framework en de Cyber Kill Chain.
De test bestond uit de volgende fasen:

\begin{enumerate}
  \item \textbf{Verkenning} -- OSINT en doelidentificatie
  \item \textbf{Bewapening} -- Tooling en payloads voorbereiden
  \item \textbf{Aflevering} -- Initiele toegang verkrijgen
  \item \textbf{Exploitatie} -- Kwetsbaarheden benutten
  \item \textbf{Installatie} -- Persistentie vestigen
  \item \textbf{Commando \& Controle} -- C2-kanaal opzetten
  \item \textbf{Acties op doel} -- Doelstellingen realiseren
\end{enumerate}
[% endif %]
[% else %]
[% if taal == 'en' %]
The penetration test was performed in accordance with the OWASP Testing Guide and PTES (Penetration Testing
Execution Standard). The test consisted of the following phases:

\begin{enumerate}
  \item \textbf{Reconnaissance} -- Gather information about the target environment
  \item \textbf{Scanning} -- Identify services and vulnerabilities
  \item \textbf{Exploitation} -- Confirm vulnerabilities
  \item \textbf{Post-exploitation} -- Determine impact
  \item \textbf{Reporting} -- Document findings
\end{enumerate}
[% else %]
De penetratietest is uitgevoerd conform de OWASP Testing Guide en PTES (Penetration Testing
Execution Standard). De test bestond uit de volgende fasen:

\begin{enumerate}
  \item \textbf{Verkenning} -- Informatie verzamelen over de doelomgeving
  \item \textbf{Scanning} -- Diensten en kwetsbaarheden identificeren
  \item \textbf{Exploitatie} -- Kwetsbaarheden bevestigen
  \item \textbf{Post-exploitatie} -- Impact bepalen
  \item \textbf{Rapportage} -- Documenteren
\end{enumerate}
[% endif %]
[% endif %]

[% for scope in testscopes %]
[% if scope == 'whitebox' %]
[% if taal == 'en' %]
The test was (partially) performed as a white box assessment. The testers had full access to source code,
architecture documentation and configurations. This allows for a thorough analysis of all security layers.
[% else %]
De test is (deels) uitgevoerd als white box assessment. De testers hadden volledige toegang tot broncode,
architectuurdocumentatie en configuraties. Dit maakt een grondige analyse mogelijk van alle
beveiligingslagen.
[% endif %]
[% elif scope == 'greybox' %]
[% if taal == 'en' %]
The test was (partially) performed as a grey box assessment. The testers had limited information, such as
network topology, user accounts or API documentation. This simulates an attacker with limited internal knowledge.
[% else %]
De test is (deels) uitgevoerd als grey box assessment. De testers hadden beperkte informatie, zoals
netwerktopologie, gebruikersaccounts of API-documentatie. Dit simuleert een aanvaller met
beperkte interne kennis.
[% endif %]
[% elif scope == 'assumebreach' %]
[% if taal == 'en' %]
The test was (partially) performed from an assume breach scenario. An already-compromised starting point within
the environment (for example a workstation or account) was assumed, so the focus is on lateral movement, privilege
escalation and detection \& response capabilities rather than on obtaining initial access.
[% else %]
De test is (deels) uitgevoerd vanuit een assume breach-scenario. Er is uitgegaan van een reeds gecompromitteerd
startpunt binnen de omgeving (bijvoorbeeld een werkplek of account), zodat de nadruk ligt op laterale beweging,
rechtenescalatie en detectie- \& responscapaciteiten in plaats van op het verkrijgen van initiële toegang.
[% endif %]
[% else %]
[% if taal == 'en' %]
The test was (partially) performed as a black box assessment. The testers had no prior information about the
internal workings, source code or architecture of the target environment. This simulates the perspective of
an external attacker.
[% else %]
De test is (deels) uitgevoerd als black box assessment. De testers hadden vooraf geen informatie over de
interne werking, broncode of architectuur van de doelomgeving. Dit simuleert het perspectief van
een externe aanvaller.
[% endif %]
[% endif %]
[% endfor %]
[% endif %]

% --- Tools en technieken ---
[% if taal == 'en' %]\subsection{Tools and Techniques}[% else %]\subsection{Tools en technieken}[% endif %]

[% if rapport_tools %]
[[ rapport_tools ]]
[% else %]
[% if taal == 'en' %]
For this assessment, use was made of \textbf{Incompetent Bastard}, an open-source project
that provides support for security assessments such as the one described in this report.
In addition, the following tools were used during the assessment:

\begin{itemize}
  \item \textbf{Nmap} -- Network discovery and port scanning
  \item \textbf{Burp Suite} -- Web application security testing
  \item \textbf{Metasploit Framework} -- Exploit development and execution
  \item \textbf{BloodHound} -- Active Directory attack path analysis
  \item \textbf{CrackMapExec} -- Network penetration testing
  \item \textbf{Impacket} -- Network protocol interaction
\end{itemize}

The specific tools used may vary depending on the scope and nature of the assessment.
All tools are used in accordance with the agreed Rules of Engagement.
[% else %]
Voor dit onderzoek is gebruik gemaakt van \textbf{Incompetent Bastard}, een opensourceproject
dat ondersteuning biedt bij beveiligingsonderzoeken zoals uitgewerkt in dit rapport.
Daarnaast zijn de volgende tools ingezet tijdens het onderzoek:

\begin{itemize}
  \item \textbf{Nmap} -- Netwerk discovery en poortscan
  \item \textbf{Burp Suite} -- Webapplicatie beveiligingstesten
  \item \textbf{Metasploit Framework} -- Exploit ontwikkeling en uitvoering
  \item \textbf{BloodHound} -- Active Directory aanvalspad-analyse
  \item \textbf{CrackMapExec} -- Netwerk penetratietesten
  \item \textbf{Impacket} -- Netwerkprotocol interactie
\end{itemize}

De specifieke tools die worden ingezet kunnen varieren afhankelijk van de scope en aard van het onderzoek.
Alle tools worden ingezet conform de overeengekomen Rules of Engagement.
[% endif %]
[% endif %]

[% endif %]
[% if sections.legsup and testtype == "redteam" and (legsup_startpunt or legsup_privileges or legsup_informatie) %]
%SEC:legsup%
% --- Legs Up ---
\subsection{Legs Up}

[% if legsup_startpunt %]
[% if taal == 'en' %]
\subsubsection{Assumed Breach Starting Points}
[% else %]
\subsubsection{Assumed Breach Startpunten}
[% endif %]

[[ legsup_startpunt | latex_escape ]]

[% endif %]
[% if legsup_privileges %]
[% if taal == 'en' %]
\subsubsection{Privilege Levels}
[% else %]
\subsubsection{Privilege Niveaus}
[% endif %]

[[ legsup_privileges | latex_escape ]]

[% endif %]
[% if legsup_informatie %]
[% if taal == 'en' %]
\subsubsection{Provided Information}
[% else %]
\subsubsection{Verstrekte Informatie}
[% endif %]

[[ legsup_informatie | latex_escape ]]

[% endif %]
[% endif %]

[# Modus operandi van de actor (alleen red team) -- vóór de scenario's. #]
[% if sections.modus_operandi and testtype == "redteam" and modus_operandi %]
%SEC:modus_operandi%
\subsection{Modus operandi}

[[ modus_operandi ]]
[% endif %]

[% if sections.scenarios and testtype == "redteam" and scenarios %]
%SEC:scenarios%
% --- Scenario's ---
[% if taal == 'en' %]
\subsection{Scenarios}
[% else %]
\subsection{Scenario's}
[% endif %]

[% for sc in scenarios %]
\subsubsection{[[ sc.naam | latex_escape ]]}

[[ sc.beschrijving | latex_escape ]]

[% endfor %]
[% endif %]

[% if sections.risicomatrix %]
%SEC:risicomatrix%
% --- Risicoinschatting ---
[% if taal == 'en' %]
\subsection{Risk Assessment}

The severity of findings is determined using the CVSS 4.0 framework
(Common Vulnerability Scoring System). The following classification is used:

\begin{center}
\begin{tabular}{|l|c|l|}
  \hline
  \rowcolor{rijkblauw}\textcolor{white}{\textbf{Classification}} & \textcolor{white}{\textbf{CVSS Score}} & \textcolor{white}{\textbf{Colour}} \\
  \hline
  Critical & 9.0 -- 10.0 & \textcolor{opaars}{Purple} \\
  High & 7.0 -- 8.9 & \textcolor{ored}{Red} \\
  Medium & 4.0 -- 6.9 & \textcolor{oorange}{Orange} \\
  Low & 0.1 -- 3.9 & \textcolor{oyellow}{Yellow} \\
  Informational & 0.0 & \textcolor{olichtblauw}{Light blue} \\
  \hline
\end{tabular}
\end{center}
[% else %]
\subsection{Risicoinschatting}

De ernst van bevindingen wordt bepaald aan de hand van het CVSS 4.0 framework
(Common Vulnerability Scoring System). De volgende classificatie wordt gehanteerd:

\begin{center}
\begin{tabular}{|l|c|l|}
  \hline
  \rowcolor{rijkblauw}\textcolor{white}{\textbf{Classificatie}} & \textcolor{white}{\textbf{CVSS Score}} & \textcolor{white}{\textbf{Kleur}} \\
  \hline
  Kritiek & 9.0 -- 10.0 & \textcolor{opaars}{Paars} \\
  Hoog & 7.0 -- 8.9 & \textcolor{ored}{Rood} \\
  Midden & 4.0 -- 6.9 & \textcolor{oorange}{Oranje} \\
  Laag & 0.1 -- 3.9 & \textcolor{oyellow}{Geel} \\
  Nihil & 0.0 & \textcolor{olichtblauw}{Lichtblauw} \\
  \hline
\end{tabular}
\end{center}
[% endif %]

[% if heatmap_rows %]
[% if taal == 'en' %]
\subsubsection*{Risk Matrix}
[% else %]
\subsubsection*{Risicomatrix}
[% endif %]

\begin{center}
\begin{tikzpicture}[x=1.25cm,y=1.25cm]
[% for row in heatmap_rows %]
[% for cell in row.cells %]
  \fill[{[[ cell.color ]]}] ([[ loop.index - 0.5 ]],[[ (row.kans|int) - 0.5 ]]) rectangle ([[ loop.index + 0.5 ]],[[ (row.kans|int) + 0.5 ]]);
  \draw[white,line width=1.2pt] ([[ loop.index - 0.5 ]],[[ (row.kans|int) - 0.5 ]]) rectangle ([[ loop.index + 0.5 ]],[[ (row.kans|int) + 0.5 ]]);
[% if cell.id_list %]  \node[font=\tiny,align=center,text width=1.15cm] at ([[ loop.index ]],[[ row.kans|int ]]) {[% for bid in cell.id_list %][% if not loop.first %], [% endif %]\hyperref[bev:[[ bid ]]]{\textcolor{white}{[[ bid ]]}}[% endfor %]};
[% endif %]
[% endfor %]
[% endfor %]
  \foreach \i in {1,...,5} { \node[below,font=\small] at (\i,0.4) {\i}; \node[left,font=\small] at (0.4,\i) {\i}; }
  \node[below=1.5em,font=\small] at (3,0.4) {Impact};
  \node[rotate=90,anchor=south,font=\small] at (-0.15,3) {[% if taal == 'en' %]Likelihood[% else %]Kans[% endif %]};
\end{tikzpicture}
\end{center}
[% endif %]
[% endif %]
[% if sections.samenvatting_bevindingen %]
%SEC:samenvatting_bevindingen%
% --- Samenvatting bevindingen ---
[% if taal == 'en' %]
\subsection{Summary of Findings}

A total of \textbf{[[ findings|length ]]} findings were identified:

\begin{center}
\begin{tabular}{|l|c|}
  \hline
  \rowcolor{rijkblauw}\textcolor{white}{\textbf{Classification}} & \textcolor{white}{\textbf{Count}} \\
  \hline
  \textcolor{opaars}{Critical} & [[ counts.critical ]] \\
  \textcolor{ored}{High} & [[ counts.high ]] \\
  \textcolor{oorange}{Medium} & [[ counts.medium ]] \\
  \textcolor{oyellow}{Low} & [[ counts.low ]] \\
  \textcolor{olichtblauw}{Informational} & [[ counts.none ]] \\
  \hline
  \textbf{Total} & \textbf{[[ findings|length ]]} \\
  \hline
\end{tabular}
\end{center}
[% else %]
\subsection{Samenvatting bevindingen}

In totaal zijn \textbf{[[ findings|length ]]} bevindingen geidentificeerd:

\begin{center}
\begin{tabular}{|l|c|}
  \hline
  \rowcolor{rijkblauw}\textcolor{white}{\textbf{Classificatie}} & \textcolor{white}{\textbf{Aantal}} \\
  \hline
  \textcolor{opaars}{Kritiek} & [[ counts.critical ]] \\
  \textcolor{ored}{Hoog} & [[ counts.high ]] \\
  \textcolor{oorange}{Midden} & [[ counts.medium ]] \\
  \textcolor{oyellow}{Laag} & [[ counts.low ]] \\
  \textcolor{olichtblauw}{Nihil} & [[ counts.none ]] \\
  \hline
  \textbf{Totaal} & \textbf{[[ findings|length ]]} \\
  \hline
\end{tabular}
\end{center}
[% endif %]

[% set _hp = sev_pie %]
[% set _hh = cvss_hist and cvss_hist.max > 0 %]
[% set _hr = remediation_bar and (remediation_bar|sum(attribute='value')) > 0 %]
[% if _hp or _hh or _hr %]
\begin{center}
% Koprij: titels op een gelijke lijn boven hun kolom
\noindent
[% if _hp -%]
\makebox[0.31\textwidth][c]{[% if taal == 'en' %]\textbf{Severity Distribution}[% else %]\textbf{Verdeling naar ernst}[% endif %]}[% if _hh or _hr %]\hfill[% endif %]
[%- endif %]
[% if _hh -%]
\makebox[0.31\textwidth][c]{[% if taal == 'en' %]\textbf{CVSS Score Distribution}[% else %]\textbf{CVSS-scoreverdeling}[% endif %]}[% if _hr %]\hfill[% endif %]
[%- endif %]
[% if _hr -%]
\makebox[0.31\textwidth][c]{[% if taal == 'en' %]\textbf{Remediation Status}[% else %]\textbf{Remediation-status}[% endif %]}
[%- endif %]
\par\vspace{6pt}
% Chartrij: afbeeldingen verticaal gecentreerd t.o.v. elkaar
\noindent
[% if _hp -%]
\begin{minipage}[c]{0.31\textwidth}\centering [[ piechart(sev_pie) ]]\end{minipage}[% if _hh or _hr %]\hfill[% endif %]
[%- endif %]
[% if _hh -%]
\begin{minipage}[c]{0.31\textwidth}\centering [[ histchart(cvss_hist, 'CVSS') ]]\end{minipage}[% if _hr %]\hfill[% endif %]
[%- endif %]
[% if _hr -%]
\begin{minipage}[c]{0.31\textwidth}\centering [[ stackbar(remediation_bar) ]]\end{minipage}
[%- endif %]
\end{center}
[% endif %]

[% endif %]
% --- Bevindingenlijst ---
[% if sections.bevindingenlijst and bevindingenlijst_rows %]
%SEC:bevindingenlijst%
[% if taal == 'en' %]
\subsection{Findings List}

The table below lists the findings identified during this assessment.
[% else %]
\subsection{Bevindingenlijst}

In de onderstaande tabel staan de bevindingen die tijdens deze test zijn aangetroffen.
[% endif %]
\begin{xltabular}{\textwidth}{|c|X|c|}
  \hline
  \rowcolor{rijkblauw}[% if taal == 'en' %]\textcolor{white}{\textbf{ID}} & \textcolor{white}{\textbf{Finding}} & \textcolor{white}{\textbf{Risk}}[% else %]\textcolor{white}{\textbf{ID}} & \textcolor{white}{\textbf{Bevinding}} & \textcolor{white}{\textbf{Risico}}[% endif %] \\
  \hline
  \endhead
\ibvervolg{3}\endfoot
\endlastfoot
[% for r in bevindingenlijst_rows %]
  \hyperref[bev:[[ r.id ]]]{[[ r.id ]]} & \hyperref[bev:[[ r.id ]]]{[[ r.titel | latex_escape ]]} & \cellcolor{[[ r.color ]]}\textbf{\textcolor{[[ r.fg ]]}{[[ r.risk ]]}} \\
  \hline
[% endfor %]
\end{xltabular}

[% endif %]
% --- Oplostijden (SLA) ---
[% if sections.sla and sla_rows %]
%SEC:sla%
[% if taal == 'en' %]
\subsubsection*{Remediation Timelines}

The recommended maximum remediation time is based on the risk classification of each finding.

\begin{xltabular}{\textwidth}{|c|X|c|p{3.5cm}|}
  \hline
  \rowcolor{rijkblauw}\textcolor{white}{\textbf{ID}} & \textcolor{white}{\textbf{Finding}} & \textcolor{white}{\textbf{Severity}} & \textcolor{white}{\textbf{Target}} \\
  \hline
  \endhead
[% else %]
\subsubsection*{Oplostijden}

De aanbevolen maximale hersteltermijn is gebaseerd op de risicoclassificatie van elke bevinding.

\begin{xltabular}{\textwidth}{|c|X|c|p{3.5cm}|}
  \hline
  \rowcolor{rijkblauw}\textcolor{white}{\textbf{ID}} & \textcolor{white}{\textbf{Bevinding}} & \textcolor{white}{\textbf{Ernst}} & \textcolor{white}{\textbf{Streeftermijn}} \\
  \hline
  \endhead
[% endif %]
\ibvervolg{4}\endfoot
\endlastfoot
[% for r in sla_rows %]
  \hyperref[bev:[[ r.id ]]]{[[ r.id ]]} & [[ r.titel | latex_escape ]] & [[ r.severity | latex_escape ]] & [[ r.sla | latex_escape ]] \\
  \hline
[% endfor %]
\end{xltabular}

\subsubsection*{[% if taal == 'en' %]Remediation Timeline[% else %]Hersteltijdlijn[% endif %]}
[[ slagantt(sla_rows) ]]
[% endif %]

% --- Detecteerbaarheid & SOC/CDC-aanbevelingen ---
[% if sections.detecteerbaarheid and detection_rated %]
%SEC:detecteerbaarheid%
[% if taal == 'en' %]
\subsubsection*{Detectability \& SOC/CDC Recommendations}
[[ hbarchart(detection_bar, '') ]]
[% if detection_gaps %]

The following findings were not (fully) detectable. Strengthening monitoring and detection for these is recommended:

\begin{itemize}
[% for g in detection_gaps %]
  \item \textbf{[[ g.id ]]: [[ g.titel | latex_escape ]]} ([[ g.label | latex_escape ]])[% if g.note %] --- [[ g.note | latex_escape ]][% endif %]
[% endfor %]
\end{itemize}
[% endif %]
[% else %]
\subsubsection*{Detecteerbaarheid \& SOC/CDC-aanbevelingen}
[[ hbarchart(detection_bar, '') ]]
[% if detection_gaps %]

De volgende bevindingen waren niet (volledig) detecteerbaar. Versterking van monitoring en detectie hiervoor wordt aanbevolen:

\begin{itemize}
[% for g in detection_gaps %]
  \item \textbf{[[ g.id ]]: [[ g.titel | latex_escape ]]} ([[ g.label | latex_escape ]])[% if g.note %] --- [[ g.note | latex_escape ]][% endif %]
[% endfor %]
\end{itemize}
[% endif %]
[% endif %]
[% endif %]

\newpage

% --- Verkenning & ontdekking (notities) ---
[% if sections.verkenning and notes %]
%SEC:verkenning%
[% if taal == 'en' %]
\section{Reconnaissance \& Discovery}
[% else %]
\section{Verkenning \& ontdekking}
[% endif %]

[% for note in notes %]
\subsection{[[ note.naam | latex_escape ]]}

[% if note.latex_raw %][[ note.uitwerken ]][% else %][[ note.uitwerken | latex_escape ]][% endif %]

[% endfor %]
\newpage
[% endif %]

[% if sections.bevindingen %]
%SEC:bevindingen%
% --- Bevindingen ---
[% if taal == 'en' %]
\section{Findings}
[% else %]
\section{Bevindingen}
[% endif %]
[# Red team-rapport: 6e metadatakolom toont MITRE ATT&CK i.p.v. NCSC #]
[% if testtype == 'redteam' %]\renewcommand{\ibnormlabel}{ATT\&CK}[% endif %]

[% for severity_label, group in grouped_findings %]
[% if group %]
[% if taal == 'en' %]
\subsection{[[ severity_label ]] findings}
[% else %]
\subsection{[[ severity_label ]] bevindingen}
[% endif %]

[% for item in group %]
\ibfindinghead{[[ '%03d' | format(item.nr) ]]}%
  {[[ item.titel | latex_escape ]]}%
  {[[ item.basescore ]]}%
  {[[ (item.owasp or '—') | latex_escape ]]}%
  {[[ (('CWE-' ~ item.cwe) if item.cwe else '—') | latex_escape ]]}%
  {[% if testtype == 'redteam' %][[ (item.mitre or '—') | latex_escape ]][% else %][[ (item.ncsc or '—') | latex_escape ]][% endif %]}

\noindent [[ item.beschrijving | latex_markup ]]
[% if item.evidence_tex %]
[[ item.evidence_tex ]]
[% endif %]
[% if item.uitwerking_detail %]

\noindent [[ item.uitwerking_detail | latex_markup ]]
[% endif %]

\ibfindingsection{Impact}
\noindent [[ item.impact | latex_markup ]]

[% if taal == 'en' %]
\ibfindingsection{Recommendation}
[% else %]
\ibfindingsection{Aanbeveling}
[% endif %]
\noindent [[ item.aanbeveling | latex_markup ]]

[% if item.locatie %]
[% if taal == 'en' %]
\noindent\textbf{Location:} \texttt{[[ item.locatie | latex_escape ]]}
[% else %]
\noindent\textbf{Locatie:} \texttt{[[ item.locatie | latex_escape ]]}
[% endif %]
\vspace{0.3em}
[% endif %]
[% if item.tester %]
[% if taal == 'en' %]
\noindent\textbf{Tested by:} [[ item.tester | latex_escape ]]
[% else %]
\noindent\textbf{Vastgesteld door:} [[ item.tester | latex_escape ]]
[% endif %]
\vspace{0.3em}
[% endif %]

[% if item.assets %]
\noindent\textbf{[% if taal == 'en' %]Affected assets:[% else %]Getroffen assets:[% endif %]} [[ item.assets | map(attribute='waarde') | join(', ') | latex_escape ]]
\vspace{0.3em}
[% endif %]

[% if item.cvss_vector %]
\noindent\textbf{CVSS Vector:} \texttt{\small [[ item.cvss_vector | latex_escape ]]}
[% endif %]

[% if item.faalmodus or item.control_ref or item.detecteerbaarheid %]
\vspace{0.3em}

\noindent
[% if item.faalmodus %][% if taal == 'en' %]\textbf{Failure mode:} [[ faalmodus_labels[item.faalmodus][1] | default(item.faalmodus) | latex_escape ]][% else %]\textbf{Faalmodus:} [[ faalmodus_labels[item.faalmodus][0] | default(item.faalmodus) | latex_escape ]][% endif %] \quad [% endif %]
[% if item.control_ref %][% if taal == 'en' %]\textbf{Control:}[% else %]\textbf{Controle:}[% endif %] [[ item.control_ref | latex_escape ]] \quad [% endif %]
[% if item.detecteerbaarheid %][% if taal == 'en' %]\textbf{Detectable:} [[ detect_labels[item.detecteerbaarheid][1] | default(item.detecteerbaarheid) | latex_escape ]][% else %]\textbf{Detecteerbaar:} [[ detect_labels[item.detecteerbaarheid][0] | default(item.detecteerbaarheid) | latex_escape ]][% endif %][% endif %]
[% if item.detectie_notitie %]

\noindent\textit{[% if taal == 'en' %]Detection notes:[% else %]Detectie-toelichting:[% endif %]} [[ item.detectie_notitie | latex_escape ]]
[% endif %]
[% endif %]

[% if item.referenties_list %]
\vspace{0.3em}
\noindent\textbf{[% if taal == 'en' %]References:[% else %]Referenties:[% endif %]}
\begin{itemize}
[% for ref in item.referenties_list %]
  \item \texttt{\small [[ ref | latex_markup ]]}
[% endfor %]
\end{itemize}
[% endif %]

[% if item.related %]
\vspace{0.3em}
\noindent\textbf{[% if taal == 'en' %]Related findings:[% else %]Gerelateerde bevindingen:[% endif %]}
\begin{itemize}
[% for rel in item.related %]
  \item [% if rel.type_label %][[ rel.type_label | latex_escape ]]: [% endif %][[ rel.naam | latex_escape ]][% if rel.nr %] (\hyperref[bev:[[ '%03d' | format(rel.nr) ]]]{[% if taal == 'en' %]finding [% else %]bevinding [% endif %][[ '%03d' | format(rel.nr) ]]})[% endif %]
[% endfor %]
\end{itemize}
[% endif %]

[% for ev in item.evidence %]
\ibevidence{[[ ev.path ]]}{[[ ev.lakpath | default('') ]]}{[[ ev.caption | latex_escape ]][% if ev.sha256 %] {\footnotesize (SHA-256: \texttt{[[ ev.sha256 ]]})}[% endif %]}
[% endfor %]

\newpage
[% endfor %]
[% endif %]
[% endfor %]

[% endif %]
[% if sections.aanbevelingen %]
%SEC:aanbevelingen%
% --- Aanbevelingen ---
[% if taal == 'en' %]
\section{Recommendations Overview}

\begin{xltabular}{\textwidth}{|c|X|c|X|}
  \hline
  \rowcolor{rijkblauw}\textcolor{white}{\textbf{ID}} & \textcolor{white}{\textbf{Finding}} & \textcolor{white}{\textbf{CVSS}} & \textcolor{white}{\textbf{Recommendation}} \\
  \hline
  \endhead
[% else %]
\section{Overzicht aanbevelingen}

\begin{xltabular}{\textwidth}{|c|X|c|X|}
  \hline
  \rowcolor{rijkblauw}\textcolor{white}{\textbf{ID}} & \textcolor{white}{\textbf{Bevinding}} & \textcolor{white}{\textbf{CVSS}} & \textcolor{white}{\textbf{Aanbeveling}} \\
  \hline
  \endhead
[% endif %]
\ibvervolg{4}\endfoot
\endlastfoot
[% for item in findings %]
  \hyperref[bev:[[ '%03d' | format(item.nr) ]]]{[[ '%03d' | format(item.nr) ]]} & [[ item.titel | latex_escape ]] & [[ item.basescore ]] & [[ (item.aanbeveling_overzicht if item.aanbeveling_overzicht else item.aanbeveling) | latex_escape ]] \\
  \hline
[% endfor %]
\end{xltabular}

[% endif %]
%SEC-END%

% --- Bijlagen (via \appendix; elk bijlage-item wordt A, B, C ...) ---
\clearpage
\appendix

[% if sections.checklists and checklists %]
% --- Checklist appendix ---
[% if taal == 'en' %]
\section{Checklist Results}
[% else %]
\section{Checklist resultaten}
[% endif %]

[% for cl in checklists %]
\subsection{[[ cl['naam'] | latex_escape ]] ([[ cl['type'] | latex_escape ]]) --- [[ cl['target'] | latex_escape ]]}

\noindent\textbf{[[ cl['stats']['pass'] ]] pass, [[ cl['stats']['fail'] ]] fail, [[ cl['stats']['warn'] ]] warn, [[ cl['stats']['skip'] ]] skip, [[ cl['stats']['na'] ]] n/a, [[ cl['stats']['open'] ]] open}
\vspace{0.5em}

[% for phase in cl['phases'] %]
\subsection{[[ phase['title'] | latex_escape ]]}

\begin{xltabular}{\textwidth}{|p{2.5cm}|X|p{3cm}|}
  \hline
  \rowcolor{rijkblauw}\textcolor{white}{\textbf{Ref}} & \textcolor{white}{\textbf{Item}} & \textcolor{white}{\textbf{Status}} \\
  \hline
  \endhead
\ibvervolg{3}\endfoot
\endlastfoot
[% for item in phase['items'] %]
  \texttt{\small [[ item['ref'] | latex_escape ]]} & [[ item['title'] | latex_escape ]] & [% if item['status'] == 'pass' %]\textcolor{ogreen}{PASS}[% elif item['status'] == 'fail' %]\textcolor{ored}{FAIL}[% elif item['status'] == 'warn' %]\textcolor{oorange}{WARN}[% elif item['status'] == 'skip' %]SKIP[% elif item['status'] == 'na' %]N/A[% else %]OPEN[% endif %] \\
  \hline
[% endfor %]
\end{xltabular}

[% endfor %]
[% endfor %]
[% endif %]

[# Blue team (detectie & respons) + Purple teaming (ART, alleen red team). #]
[% if sections.blueteam and testtype == "redteam" and blueteam %]
%SEC:blueteam%
\section{[% if taal == 'en' %]Blue team --- detection \& response[% else %]Blue team --- detectie \& respons[% endif %]}

[[ blueteam ]]
[% endif %]
[% if sections.purpleteam and testtype == "redteam" and purpleteam %]
%SEC:purpleteam%
\section{Purple teaming}

[[ purpleteam ]]
[% endif %]

[% if (owasp_rows or ncsc_rows or controls_rows or chain_paths) and not checklists %]
[% endif %]

[% if sections.killchain and chain_paths %]
[% if taal == 'en' %]
\section{Exploit Chains / Kill Chain}
[% else %]
\section{Exploit chains / kill chain}
[% endif %]
[% for cp in chain_paths %]
\noindent\textbf{[% if taal == 'en' %]Chain[% else %]Keten[% endif %] [[ cp.num ]]}\par\nopagebreak
\begin{center}
\begin{tikzpicture}[node distance=5mm and 0mm]
[% for n in cp.nodes %]
  \node[draw,rounded corners,fill=cblue!12,inner sep=5pt,text width=0.72\textwidth,align=left,font=\small[% if not loop.first %],below=of n[[ loop.index0 ]][% endif %]] (n[[ loop.index ]]) {\textbf{[[ n.id ]]}~[[ n.naam | latex_escape ]]};
[% endfor %]
[% for n in cp.nodes %][% if not loop.first %]
  \draw[-{Stealth[length=2mm]},thick] (n[[ loop.index0 ]]) -- (n[[ loop.index ]]);
[% endif %][% endfor %]
\end{tikzpicture}
\end{center}
[% endfor %]
[% endif %]

[% if sections.owasp_mapping and owasp_rows %]
[% if taal == 'en' %]
\section{OWASP Top 10 Mapping}
\begin{xltabular}{\textwidth}{|c|X|X|}
  \hline
  \rowcolor{rijkblauw}\textcolor{white}{\textbf{ID}} & \textcolor{white}{\textbf{Finding}} & \textcolor{white}{\textbf{OWASP 2025}} \\
  \hline
  \endhead
[% else %]
\section{OWASP Top 10 mapping}
\begin{xltabular}{\textwidth}{|c|X|X|}
  \hline
  \rowcolor{rijkblauw}\textcolor{white}{\textbf{ID}} & \textcolor{white}{\textbf{Bevinding}} & \textcolor{white}{\textbf{OWASP 2025}} \\
  \hline
  \endhead
[% endif %]
\ibvervolg{3}\endfoot
\endlastfoot
[% for r in owasp_rows %]
  \hyperref[bev:[[ r.id ]]]{[[ r.id ]]} & [[ r.titel | latex_escape ]] & [% if r.owasp_2025_url %]\href{[[ r.owasp_2025_url ]]}{[[ r.owasp_2025 | latex_escape ]]}[% else %][[ r.owasp_2025 | latex_escape ]][% endif %] \\
  \hline
[% endfor %]
\end{xltabular}
[% endif %]

[# Red team-rapport: ATT&CK-mapping i.p.v. NCSC-richtlijnen mapping #]
[% if sections.ncsc_mapping and testtype == 'redteam' and mitre_rows %]
\section{[% if taal == 'en' %]MITRE ATT\&CK Mapping[% else %]MITRE ATT\&CK-mapping[% endif %]}
\begin{xltabular}{\textwidth}{|c|X|X|}
  \hline
  \rowcolor{rijkblauw}\textcolor{white}{\textbf{ID}} & \textcolor{white}{\textbf{[% if taal == 'en' %]Finding[% else %]Bevinding[% endif %]}} & \textcolor{white}{\textbf{ATT\&CK}} \\
  \hline
  \endhead
\ibvervolg{3}\endfoot
\endlastfoot
[% for r in mitre_rows %]
  \hyperref[bev:[[ r.id ]]]{[[ r.id ]]} & [[ r.titel | latex_escape ]] & [[ r.mitre | latex_escape ]] \\
  \hline
[% endfor %]
\end{xltabular}
[% elif sections.ncsc_mapping and ncsc_rows %]
[% if taal == 'en' %]
\section{NCSC Guidelines Mapping}
\begin{xltabular}{\textwidth}{|c|X|X|}
  \hline
  \rowcolor{rijkblauw}\textcolor{white}{\textbf{ID}} & \textcolor{white}{\textbf{Finding}} & \textcolor{white}{\textbf{NCSC}} \\
  \hline
  \endhead
[% else %]
\section{NCSC-richtlijnen mapping}
\begin{xltabular}{\textwidth}{|c|X|X|}
  \hline
  \rowcolor{rijkblauw}\textcolor{white}{\textbf{ID}} & \textcolor{white}{\textbf{Bevinding}} & \textcolor{white}{\textbf{NCSC}} \\
  \hline
  \endhead
[% endif %]
\ibvervolg{3}\endfoot
\endlastfoot
[% for r in ncsc_rows %]
  \hyperref[bev:[[ r.id ]]]{[[ r.id ]]} & [[ r.titel | latex_escape ]] & [[ r.ncsc | latex_escape ]] \\
  \hline
[% endfor %]
\end{xltabular}
[% endif %]

[% if sections.controls_mapping and controls_rows %]
[% if taal == 'en' %]
\section{ISO 27001 / BIO / NIS2 Mapping}
\footnotesize\emph{Indicative mapping derived from the finding CWE via OWASP Top 10:2025 category; not a formal certification or complete coverage.}\normalsize
\begin{xltabular}{\textwidth}{|c|X|X|X|X|}
  \hline
  \rowcolor{rijkblauw}\textcolor{white}{\textbf{ID}} & \textcolor{white}{\textbf{Finding}} & \textcolor{white}{\textbf{ISO 27001:2022}} & \textcolor{white}{\textbf{BIO2}} & \textcolor{white}{\textbf{NIS2 art. 21}} \\
  \hline
  \endhead
[% else %]
\section{ISO 27001 / BIO / NIS2 mapping}
\footnotesize\emph{Indicatieve koppeling, afgeleid van de CWE via de OWASP Top 10:2025-categorie; geen formele certificering of volledige dekking.}\normalsize
\begin{xltabular}{\textwidth}{|c|X|X|X|X|}
  \hline
  \rowcolor{rijkblauw}\textcolor{white}{\textbf{ID}} & \textcolor{white}{\textbf{Bevinding}} & \textcolor{white}{\textbf{ISO 27001:2022}} & \textcolor{white}{\textbf{BIO2}} & \textcolor{white}{\textbf{NIS2 art. 21}} \\
  \hline
  \endhead
[% endif %]
\ibvervolg{5}\endfoot
\endlastfoot
[% for r in controls_rows %]
  \hyperref[bev:[[ r.id ]]]{[[ r.id ]]} & [[ r.titel | latex_escape ]] & [[ r.iso27001 | latex_escape ]] & [[ r.bio | latex_escape ]] & [[ r.nis2 | latex_escape ]] \\
  \hline
[% endfor %]
\end{xltabular}
[% endif %]

[% if sections.cia_triad and cia_triad %]
[% if taal == 'en' %]\section{CIA Triad}[% else %]\section{CIA-triade}[% endif %]
\footnotesize\emph{[[ cia_triad.intro | latex_escape ]]}\normalsize
\begin{center}
\begin{tikzpicture}[scale=1]
  \coordinate (C) at (90:3);
  \coordinate (I) at (330:3);
  \coordinate (A) at (210:3);
  \draw[gray!45] (C) -- (I) -- (A) -- cycle;
  \draw[gray!25] (90:2) -- (330:2) -- (210:2) -- cycle;
  \draw[gray!25] (90:1) -- (330:1) -- (210:1) -- cycle;
  \draw[gray!45] (0,0) -- (C) (0,0) -- (I) (0,0) -- (A);
  \fill[blue!35,opacity=0.55] (90:[[ cia_triad.c_r ]]) -- (330:[[ cia_triad.i_r ]]) -- (210:[[ cia_triad.a_r ]]) -- cycle;
  \draw[blue!70,thick] (90:[[ cia_triad.c_r ]]) -- (330:[[ cia_triad.i_r ]]) -- (210:[[ cia_triad.a_r ]]) -- cycle;
  \node[above] at (C) {[% if taal == 'en' %]Confidentiality[% else %]Vertrouwelijkheid[% endif %] ([[ cia_triad.c ]]\%)};
  \node[below right] at (I) {[% if taal == 'en' %]Integrity[% else %]Integriteit[% endif %] ([[ cia_triad.i ]]\%)};
  \node[below left] at (A) {[% if taal == 'en' %]Availability[% else %]Beschikbaarheid[% endif %] ([[ cia_triad.a ]]\%)};
\end{tikzpicture}
\end{center}
[% endif %]

[% if sections.subsequent_triad and subsequent_triad %]
[% if taal == 'en' %]\section{Subsequent System Impact}[% else %]\section{Vervolgsysteem-impact}[% endif %]
\footnotesize\emph{[[ subsequent_triad.intro | latex_escape ]]}\normalsize
\begin{center}
\begin{tikzpicture}[scale=1]
  \coordinate (C) at (90:3);
  \coordinate (I) at (330:3);
  \coordinate (A) at (210:3);
  \draw[gray!45] (C) -- (I) -- (A) -- cycle;
  \draw[gray!25] (90:2) -- (330:2) -- (210:2) -- cycle;
  \draw[gray!25] (90:1) -- (330:1) -- (210:1) -- cycle;
  \draw[gray!45] (0,0) -- (C) (0,0) -- (I) (0,0) -- (A);
  \fill[orange!35,opacity=0.55] (90:[[ subsequent_triad.c_r ]]) -- (330:[[ subsequent_triad.i_r ]]) -- (210:[[ subsequent_triad.a_r ]]) -- cycle;
  \draw[orange!80,thick] (90:[[ subsequent_triad.c_r ]]) -- (330:[[ subsequent_triad.i_r ]]) -- (210:[[ subsequent_triad.a_r ]]) -- cycle;
  \node[above] at (C) {[% if taal == 'en' %]Confidentiality[% else %]Vertrouwelijkheid[% endif %] ([[ subsequent_triad.c ]]\%)};
  \node[below right] at (I) {[% if taal == 'en' %]Integrity[% else %]Integriteit[% endif %] ([[ subsequent_triad.i ]]\%)};
  \node[below left] at (A) {[% if taal == 'en' %]Availability[% else %]Beschikbaarheid[% endif %] ([[ subsequent_triad.a ]]\%)};
\end{tikzpicture}
\end{center}
[% endif %]

[% if sections.exploit_profile and exploit_profile %]
[% if taal == 'en' %]\section{Exploitability Profile}[% else %]\section{Exploiteerbaarheids-profiel}[% endif %]
\footnotesize\emph{[% if taal == 'en' %]Distribution of CVSS exploitability metrics across [[ exploit_profile.n ]] findings.[% else %]Verdeling van de CVSS-exploiteerbaarheidsmetrieken over [[ exploit_profile.n ]] bevindingen.[% endif %]}\normalsize
[% macro exptable(metrics, taal) %]
\begin{tabularx}{\linewidth}{|l|X|c|}
  \hline
  \rowcolor{rijkblauw}[% if taal == 'en' %]\textcolor{white}{\textbf{Metric}} & \textcolor{white}{\textbf{Value}} & \textcolor{white}{\textbf{Count}}[% else %]\textcolor{white}{\textbf{Metriek}} & \textcolor{white}{\textbf{Waarde}} & \textcolor{white}{\textbf{Aantal}}[% endif %] \\
  \hline
[% for met in metrics %]
[% for it in met['items'] %]
  [% if loop.first %][[ met['name'] | latex_escape ]][% endif %] & [[ it['label'] | latex_escape ]] & [[ it['count'] ]] \\
  \hline
[% endfor %]
[% endfor %]
\end{tabularx}
[% endmacro %]
[% if exploit_profile['gauge_angle'] %]
\noindent
\begin{minipage}[c]{0.44\textwidth}
\ibgauge{[[ exploit_profile['gauge_angle'] ]]}{[[ exploit_profile['index'] ]]\%}{[[ exploit_profile['index_label'] | latex_escape ]]}{[% if taal == 'en' %]Exploitability index (indicative)[% else %]Exploiteerbaarheidsindex (indicatief)[% endif %]}
\end{minipage}\hfill
\begin{minipage}[c]{0.53\textwidth}
[[ exptable(exploit_profile['metrics'], taal) ]]
\end{minipage}
[% else %]
[[ exptable(exploit_profile['metrics'], taal) ]]
[% endif %]
[% endif %]

[% if sections.direct_exploitable and direct_rows %]
[% if taal == 'en' %]\section{Directly Exploitable Findings}[% else %]\section{Direct exploiteerbare bevindingen}[% endif %]
\footnotesize\emph{[[ direct_intro | latex_escape ]]}\normalsize
\begin{xltabular}{\textwidth}{|c|X|c|}
  \hline
  \rowcolor{rijkblauw}[% if taal == 'en' %]\textcolor{white}{\textbf{ID}} & \textcolor{white}{\textbf{Finding}} & \textcolor{white}{\textbf{Risk}}[% else %]\textcolor{white}{\textbf{ID}} & \textcolor{white}{\textbf{Bevinding}} & \textcolor{white}{\textbf{Risico}}[% endif %] \\
  \hline
  \endhead
\ibvervolg{3}\endfoot
\endlastfoot
[% for r in direct_rows %]
  \hyperref[bev:[[ r.id ]]]{[[ r.id ]]} & \hyperref[bev:[[ r.id ]]]{[[ r.titel | latex_escape ]]} & \cellcolor{[[ r.color ]]}\textbf{\textcolor{[[ r.fg ]]}{[[ r.risk ]]}} \\
  \hline
[% endfor %]
\end{xltabular}
[% endif %]

[% if sections.impact_gauge and impact_gauge %]
[% if taal == 'en' %]\section{Impact Meter}[% else %]\section{Impact-meter}[% endif %]
\footnotesize\emph{[% if taal == 'en' %]Aggregate impact index across [[ impact_gauge['n'] ]] findings, from the CVSS C/I/A impact metrics (indicative).[% else %]Geaggregeerde impact-index over [[ impact_gauge['n'] ]] bevindingen, afgeleid van de CVSS C/I/A-impactmetrieken (indicatief).[% endif %]}\normalsize
\ibgauge{[[ impact_gauge['gauge_angle'] ]]}{[[ impact_gauge['index'] ]]\%}{[[ impact_gauge['index_label'] | latex_escape ]]}{[% if taal == 'en' %]Impact index (indicative)[% else %]Impact-index (indicatief)[% endif %]}
[% endif %]

[% if sections.retest_voortgang and retest_progress %]
[% if taal == 'en' %]\section{Retest Progress}[% else %]\section{Retest-voortgang}[% endif %]
\footnotesize\emph{[[ retest_intro ]]}\normalsize
\par\vspace{4pt}
[[ stackbar(retest_progress) ]]
\par\vspace{4pt}
\begin{tabular}{@{}ll@{}}
[% for r in retest_progress %]
\textcolor{[[ r.color ]]}{\rule{1.4ex}{1.4ex}}~[[ r.label ]] & [[ r.value ]] \\
[% endfor %]
\end{tabular}
[% if retest_open %]
\par\vspace{4pt}\noindent\textbf{[% if taal == 'en' %]Still open after retest:[% else %]Nog open na retest:[% endif %]}
\begin{itemize}
[% for o in retest_open %]
  \item \hyperref[bev:[[ o.id ]]]{[[ o.id ]]} --- [[ o.titel ]]
[% endfor %]
\end{itemize}
[% endif %]
[% endif %]

[% if sections.beperkingen %]
[% if taal == 'en' %]\section{Scope Limitations \& Assumptions}[% else %]\section{Beperkingen en aannames}[% endif %]
\begin{itemize}
[% for ln in beperkingen_lines %]
  \item [[ ln ]]
[% endfor %]
\end{itemize}
[% endif %]

[% if sections.cvss_detail and cvss_detail_rows %]
[% if taal == 'en' %]\section{CVSS Vector Breakdown}[% else %]\section{CVSS-vectordetail}[% endif %]
\footnotesize\emph{[% if taal == 'en' %]Decoding of the CVSS base vector per finding, so the assigned score is transparent and auditable.[% else %]Ontleding van de CVSS-base-vector per bevinding, zodat de toegekende score transparant en auditbaar is.[% endif %]}\normalsize
\begin{xltabular}{\textwidth}{|c|l|c|>{\raggedright\arraybackslash}X|}
  \hline
  \rowcolor{rijkblauw}[% if taal == 'en' %]\textcolor{white}{\textbf{ID}} & \textcolor{white}{\textbf{Finding}} & \textcolor{white}{\textbf{CVSS}} & \textcolor{white}{\textbf{Metrics}}[% else %]\textcolor{white}{\textbf{ID}} & \textcolor{white}{\textbf{Bevinding}} & \textcolor{white}{\textbf{CVSS}} & \textcolor{white}{\textbf{Metrieken}}[% endif %] \\
  \hline
  \endhead
\ibvervolg{4}\endfoot
\endlastfoot
[% for r in cvss_detail_rows %]
  \hyperref[bev:[[ r.id ]]]{[[ r.id ]]} & \hyperref[bev:[[ r.id ]]]{[[ r.titel ]]} & [[ r.score ]] & {\footnotesize [[ r.summary ]]} \\
  \hline
[% endfor %]
\end{xltabular}
[% endif %]

[% if sections.bewijsregister and evidence_register %]
[% if taal == 'en' %]\section{Evidence Register}[% else %]\section{Bewijsregister}[% endif %]
\footnotesize\emph{[% if taal == 'en' %]Each evidence file is listed with its SHA-256 hash so that its integrity can be independently verified.[% else %]Elk bewijsbestand is opgenomen met zijn SHA-256-hash, zodat de integriteit onafhankelijk geverifieerd kan worden.[% endif %]}\normalsize
\begin{xltabular}{\textwidth}{|c|>{\raggedright\arraybackslash}X|l|}
  \hline
  \rowcolor{rijkblauw}[% if taal == 'en' %]\textcolor{white}{\textbf{ID}} & \textcolor{white}{\textbf{Evidence}} & \textcolor{white}{\textbf{SHA-256}}[% else %]\textcolor{white}{\textbf{ID}} & \textcolor{white}{\textbf{Bewijs}} & \textcolor{white}{\textbf{SHA-256}}[% endif %] \\
  \hline
  \endhead
\ibvervolg{3}\endfoot
\endlastfoot
[% for r in evidence_register %]
  \hyperref[bev:[[ r.id ]]]{[[ r.id ]]} & [[ r.caption ]][% if r.file %] {\footnotesize\itshape([[ r.file ]])}[% endif %] & {\ttfamily\scriptsize [[ r.sha_fmt ]]} \\
  \hline
[% endfor %]
\end{xltabular}
[% endif %]

[% if sections.begrippenlijst %]
[% if taal == 'en' %]\section{Glossary \& Abbreviations}[% else %]\section{Begrippen- en afkortingenlijst}[% endif %]
\begin{xltabular}{\textwidth}{|l|>{\raggedright\arraybackslash}X|}
  \hline
  \rowcolor{rijkblauw}[% if taal == 'en' %]\textcolor{white}{\textbf{Term}} & \textcolor{white}{\textbf{Description}}[% else %]\textcolor{white}{\textbf{Begrip}} & \textcolor{white}{\textbf{Omschrijving}}[% endif %] \\
  \hline
  \endhead
\ibvervolg{2}\endfoot
\endlastfoot
[% for r in begrippen_rows %]
  \textbf{[[ r.term ]]} & [[ r.definitie ]] \\
  \hline
[% endfor %]
\end{xltabular}
[% endif %]

[# Verklaring van betrokken partijen (ART, alleen red team). #]
[% if sections.verklaring and testtype == "redteam" and verklaring %]
%SEC:verklaring%
\section{[% if taal == 'en' %]Declaration[% else %]Verklaring[% endif %]}

[[ verklaring ]]
[% endif %]

[# Eigen vrije-tekst bijlagen (door de gebruiker toegevoegd, na \appendix). #]
[% for b in bijlagen %]
\section{[[ b.titel ]]}

[[ b.tekst ]]
[% endfor %]

\end{document}
